

\section{Introduction}
\label{sec:intro}

  Modeling and learning from sequential data is a fundamental problem in
  modern machine learning, underlying tasks such as language modeling, speech
  recognition, video processing, and reinforcement learning.
  A core aspect of modeling long-term and complex temporal dependencies is \emph{memory}, or
  storing and incorporating information from previous time steps.
  The challenge is learning a representation of the entire cumulative history using bounded storage,
  which must be updated online as more data is received.
  

  One established approach is to model
  a state that evolves over time as it incorporates more information.
  The deep learning instantiation of this approach is the recurrent neural network (RNN),
  which is known to suffer from a limited memory horizon~\citep{lstm,jaeger2004harnessing,pascanu2013difficulty} (e.g., the ``vanishing gradients'' problem).
  Although various heuristics have been proposed to overcome this,
  such as gates in the successful LSTM and GRU~\citep{lstm, cho2014learning},
  or higher-order frequencies in the recent Fourier Recurrent
  Unit~\citep{zhang2018learning} and Legendre Memory Unit
  (LMU)~\citep{voelker2019legendre},
  a unified understanding of memory remains a challenge.
  Furthermore, existing methods generally require priors on the sequence length or timescale and are ineffective outside this range~\citep{tallec2018can, voelker2019legendre}; this can be problematic in settings with distribution shift (e.g.\ arising from different instrument sampling rates in medical data~\cite{saab2020weak,shah2018temple}).
  Finally, many of them lack theoretical guarantees on how well they capture long-term
  dependencies, such as gradient bounds.
  To design a better memory representation, we would ideally
  (i) have a unified view of these existing methods,
  (ii) be able to address dependencies of any length without priors on the timescale,
  and
  (iii) have a rigorous theoretical understanding of their memory mechanism.







  Our insight is to phrase \emph{memory} as a technical problem of \emph{online
    function approximation} where a function $f(t) : \R_+ \to \R$ is summarized by storing its
  optimal coefficients in terms of some basis functions.
  This approximation is evaluated with respect to a measure that
  specifies the importance of each time in the past.
  Given this function approximation formulation, orthogonal polynomials (OPs) emerge as a natural basis since
  their optimal coefficients can be expressed in closed form~\citep{chihara}.
  With their rich and well-studied history~\citep{szego}, along with their
  widespread use in approximation theory~\citep{trefethen2019approximation} and
  signal processing~\citep{proakis2001digital}, OPs bring a library of
  techniques to this memory representation problem.
  We formalize a framework, \textbf{HiPPO} (high-order polynomial projection
  operators), which produces operators that project arbitrary functions onto the space of
  orthogonal polynomials with respect to a given measure.
  This general framework allows us to analyze several families of measures,
  where this operator, as a closed-form ODE or linear recurrence, allows fast incremental updating of the optimal polynomial approximation as the input function is revealed through time.


  By posing a formal optimization problem underlying recurrent sequence models,
  the HiPPO framework (Section~\ref{sec:framework})
  generalizes and explains previous methods, unlocks new methods
  appropriate for sequential data at different timescales, and comes with several theoretical guarantees.
  (i) For example, with a short derivation we exactly recover as a special case
  the LMU~\citep{voelker2019legendre} (Section~\ref{subsec:high_order_projection}), which proposes an update rule that projects
  onto fixed-length sliding windows through time.%
  \footnote{The LMU was originally motivated by spiking neural networks in
    modeling biological nervous systems; its derivation is not self-contained but a
    sketch can be pieced together from~\citep{voelker2019legendre,voelker2018improving,voelker2019dynamical}.}
  HiPPO also sheds new light on classic techniques such as the gating mechanism
  of LSTMs and GRUs, which arise in one extreme using only low-order degrees in
  the approximation (Section~\ref{subsec:low_order_projection}).
  (ii)
    By choosing more suitable measures, HiPPO yields a novel mechanism (Scaled Legendre, or LegS) that
    always takes into account the function's full history instead of a sliding window.
    This flexibility removes the need for hyperparameters or priors on the sequence length,
    allowing LegS to generalize to different input timescales.
    (iii)
    The connections to dynamical systems and approximation theory
    allows us to show several theoretical benefits of HiPPO-LegS:
    invariance to input timescale, asymptotically more efficient updates, and
    bounds on gradient flow and approximation error (Section~\ref{sec:theory_legs}).
    

  We integrate the HiPPO memory mechanisms into RNNs, and empirically show that they outperform baselines on standard tasks used to benchmark long-term dependencies.
  On the permuted MNIST dataset, %
  our hyperparameter-free HiPPO-LegS method achieves a new state-of-the-art
  accuracy of 98.3\%,
  beating the previous RNN SoTA by over 1 point and even outperforming models
  with global context such as transformers (Section~\ref{subsec:membenchmark}).
  Next, we demonstrate the timescale robustness of HiPPO-LegS on a novel
  trajectory classification task, where it is able to generalize to unseen
  timescales and handle missing data whereas RNN and neural ODE baselines fail
  (Section~\ref{subsec:exp-timescale}).
  Finally, we validate HiPPO's theory, including computational efficiency and
  scalability, allowing fast and accurate online function reconstruction over
  millions of time steps (Section~\ref{subsec:exp-scalability}).
  Code for reproducing our experiments is available at \url{https://github.com/HazyResearch/hippo-code}.







